\documentclass[10pt,a4paper]{amsart}
\usepackage[margin=19mm]{geometry}
\usepackage{amsmath,amssymb,mathtools}
\usepackage[scr=rsfs,cal=euler]{mathalfa}
\usepackage{xfrac}
\usepackage[unicode,hidelinks,bookmarks=false]{hyperref}
\newcommand{\ie}{i.e.\ }
\newcommand{\cf}{cf.\ }
\newcommand{\resp}{respectively\ }
\newcommand{\Subsection}{\subsection}
\providecommand{\nobreakdash}{\nobreak\hbox{-}}
\setlength{\emergencystretch}{2em}
\multlinegap0pt
\usepackage{amssymb}
\usepackage[shortlabels]{enumitem}
\newtheorem{lemma}{Lemma}
\newcommand{\C}{\mathbb C}
\newcommand{\R}{\mathbb R}
\newcommand{\T}{\mathbb T}
\newcommand{\ol}{\overline}
\newcommand{\pd}{\partial}
\title{Geometric Orbital Linearization of Siegel Singularities in the Plane}
\author{Toshikazu Ito \and Bruno Sc\'ardua}
\date{Source: arXiv:2609.14855v1 (CC BY 4.0)}
\begin{document}
\maketitle

\noindent\emph{Source and licence.} Geometric Orbital Linearization of Siegel Singularities in the Plane. Version arXiv:2609.14855v1 (CC BY 4.0), distributed by arXiv under the Creative Commons Attribution 4.0 International licence, http://creativecommons.org/licenses/by/4.0/, as recorded in the arXiv licence metadata for this exact version and shown on the abstract page https://arxiv.org/abs/2609.14855v1. Full licence notice: \texttt{licenses/arxiv-2609.14855-CC-BY-4.0.txt}.\par\smallskip

\noindent\textit{Editorial scope.} Counted unit: the article's lemma lem:reinhardt-rigidity (source0.tex lines 485--503). The article's complete original proof of it is delivered with it (source0.tex lines 505--639, one \texttt{proof} environment of 135 lines). No mathematics is written, repaired or re-derived: the statement and the proof are the article's own lines, in the article's own notation, and no retained line is substituted or re-flowed.  Retained uncounted as the source-local context the proof needs: lines 430--432; lines 439--445.  Nothing was omitted from any retained span, so the package carries no omission record.  No retained line contains a citation, so the wrapper carries no bibliography and no bibliography entry.  No retained line contains a figure environment, an image-inclusion command or drawing code, so no image is delivered and the package declares no asset dependency.  Editorial formatting only, in addition: an AMS \texttt{amsart} preamble that restates the article's own declarations and macros used by the retained lines, namely the environments \texttt{lemma} on separate counters, together with the standard packages those lines need.  The retained environments print in this document as Lemma 1 in order of appearance, whereas the article prints the same items with the numbering of its own full document.  The complete native source of the article as distributed in the arXiv e-print for this exact version is bundled unchanged as \texttt{sources/210367/source0.tex}.
\begin{equation}\label{eq:reinhardt-profile}
 t=\varphi(s),\qquad 0<s<a,
\end{equation}

\begin{lemma}[Levi monotonicity of the Reinhardt profile]\label{lem:reinhardt-profile}
With the notation above, the function
\[
 h(s)=-\frac{s\varphi'(s)}{\varphi(s)},\qquad 0<s<a,
\]
is strictly increasing and maps $(0,a)$ onto $(0,+\infty)$.
\end{lemma}

\begin{lemma}[fixed-coordinate Reinhardt rigidity]\label{lem:reinhardt-rigidity}
Let $X$ be a holomorphic vector-field germ at $0\in\C^2$ with isolated nondegenerate singularity, and suppose that in the chosen coordinates its linear part is diagonal,
\[
 X_1=\alpha z\frac{\pd}{\pd z}+\beta w\frac{\pd}{\pd w},
 \qquad \alpha\beta\ne0.
\]
Assume that there exists a sequence $r_j\downarrow0$ such that, for every $j$, the complex tangency locus $M(X,\partial(r_jD))$ is a nonempty smooth real surface of dimension two. Then
\[
 \frac{\alpha}{\beta}\in\R_{<0},
\]
and, after multiplication of $X$ by a nonzero constant,
\[
 X=h_0(z,w)\left(
 \lambda z\frac{\pd}{\pd z}-\mu w\frac{\pd}{\pd w}
 \right),
 \qquad \lambda,\mu>0,
\]
for a holomorphic unit $h_0$.
\end{lemma}

\begin{proof}
Rescale $z=r\zeta$, $w=r\omega$. The vector field
\[
 \widetilde X_r=r^{-1}X(r\zeta,r\omega)
\]
converges in $C^1$ on compact sets to $X_1$. Choose a tangency point on each $\partial(r_jD)$ and rescale it to $p_j\in\partial D$. Passing to a subsequence, $p_j\to p_0\in\partial D$, and
\[
 X_1(\rho)(p_0)=0
\]
for any local Reinhardt defining function $\rho$.

The point $p_0$ cannot lie on a coordinate axis. For example, at a point $(0,w)$ of $\partial D$ with $w\ne0$, Reinhardt invariance gives $\rho_z=0$, while regularity of the boundary gives $\rho_w\ne0$; hence $X_1(\rho)=\beta w\rho_w\ne0$. The argument on the other axis is identical. Thus $p_0$ has both coordinates nonzero. Writing locally $\rho=\rho(s,t)$, the complete Reinhardt property together with strict pseudoconvexity (equivalently, with the monotonicity of the profile established in Lemma~\ref{lem:reinhardt-profile}) gives $\rho_s,\rho_t>0$ along the positive part of the boundary, after choosing the outward orientation of $\rho$, and therefore
\[
 \alpha s\rho_s+\beta t\rho_t=0
\]
at $p_0$. It follows that
\[
 \frac{\alpha}{\beta}=-\frac{t\rho_t}{s\rho_s}\in\R_{<0}.
\]
After multiplication by a nonzero constant we may write
\[
 X_1=X_0=\lambda z\frac{\pd}{\pd z}-\mu w\frac{\pd}{\pd w},
 \qquad \lambda,\mu>0.
\]

For the profile~\eqref{eq:reinhardt-profile},
\begin{equation}\label{eq:linear-reinhardt-tangency}
 X_0(\rho)
 =-\lambda s\varphi'(s)-\mu\varphi(s)
 =\varphi(s)\bigl(\lambda h(s)-\mu\bigr).
\end{equation}
By Lemma~\ref{lem:reinhardt-profile}, this vanishes on exactly one Reinhardt torus
\[
 T_D=\{|z|^2=s_0,\ |w|^2=t_0\},
 \qquad t_0=\varphi(s_0),
\]
where $h(s_0)=\mu/\lambda$. The zero is transverse, since
\begin{equation}\label{eq:transverse-reinhardt-zero}
 \frac{d}{ds}X_0(\rho)\Big|_{s=s_0}
 =\lambda\varphi(s_0)h'(s_0)>0.
\end{equation}

Suppose now, toward a contradiction, that $X\wedge X_0\not\equiv0$. Multiplication by a holomorphic unit does not change the complex tangency set. We may therefore remove, one degree at a time, every homogeneous term which is a holomorphic multiple of $X_0$ before the first degree at which the wedge with $X_0$ is nonzero. Since only finitely many degrees precede that first degree, this requires only finitely many unit multiplications. Thus, for some $m\ge2$,
\[
 X=X_0+X_m+O(m+1),
 \qquad
 X_m=P_m(z,w)\frac{\pd}{\pd z}+Q_m(z,w)\frac{\pd}{\pd w},
\]
where $P_m,Q_m$ are homogeneous of degree $m$ and $X_m$ is not a homogeneous holomorphic multiple of $X_0$.

On the fixed boundary $\partial D$, after rescaling, put
\[
 G_r=\widetilde X_r(\rho).
\]
Then
\begin{equation}\label{eq:Gr-reinhardt}
 G_r=G_0+r^{m-1}F_m+O(r^m),
\end{equation}
where
\[
 G_0=X_0(\rho)
\]
and
\[
 F_m=\rho_sP_m\ol z+\rho_tQ_m\ol w.
\]
The convergence is uniform with first derivatives on $\partial D$.

By~\eqref{eq:linear-reinhardt-tangency}--\eqref{eq:transverse-reinhardt-zero}, $G_0$ has exactly one zero torus $T_D$ and is transverse to zero there. Choose a small Reinhardt tubular neighborhood $U$ of $T_D$. On the compact set $\partial D\setminus U$, $|G_0|$ is bounded below by a positive constant, so for sufficiently small $r$ the equation $\Re G_r=0$ has no solutions there. Inside $U$, the implicit-function theorem and the uniform nonvanishing radial derivative show that $\Re G_r=0$ is a unique connected global graph
\[
 \Sigma_r\simeq\T^2
\]
over $T_D$.

For $r=r_j$, the tangency locus
\[
 M_r=\{G_r=0\}
\]
is, by hypothesis, a nonempty smooth compact real surface of dimension two contained in the connected two-manifold $\Sigma_r$. Its inclusion is locally open, while compactness makes it closed. Hence
\[
 M_r=\Sigma_r.
\]
In particular, $\Im G_r=0$ on the whole graph. Dividing the imaginary part of~\eqref{eq:Gr-reinhardt} by $r_j^{m-1}$ and letting $j\to\infty$ gives
\[
 \Im F_m=0
\]
on $T_D$.

On $T_D$ the quantities
\[
 c_1=\rho_s(s_0,t_0),\qquad c_2=\rho_t(s_0,t_0)
\]
are positive constants, so
\[
 F_m=c_1P_m\ol z+c_2Q_m\ol w.
\]
Every Fourier frequency occurring in $P_m\ol z$ or $Q_m\ol w$ has component sum $m-1>0$. A real-valued function on $\T^2$ must contain with each frequency its opposite conjugate frequency, whose component sum would be $-(m-1)$. No such frequencies occur. Therefore
\[
 F_m\equiv0\qquad\text{on }T_D.
\]
Since on $T_D$
\[
 \ol z=\frac{s_0}{z},\qquad \ol w=\frac{t_0}{w},
\]
we obtain
\[
 c_1s_0\,wP_m+c_2t_0\,zQ_m=0
\]
on $T_D$, hence identically as a holomorphic polynomial. The linear tangency relation on $T_D$ is
\[
 \lambda c_1s_0=\mu c_2t_0,
\]
and consequently
\[
 \mu wP_m+\lambda zQ_m\equiv0.
\]
Since $z$ and $w$ are relatively prime,
\[
 P_m=zR_{m-1},\qquad
 Q_m=-\frac{\mu}{\lambda}wR_{m-1}
\]
for a homogeneous polynomial $R_{m-1}$. Thus $X_m=(R_{m-1}/\lambda)X_0$, contradicting the choice of $m$. Hence
\[
 X\wedge X_0\equiv0.
\]
Writing $X=P\partial_z+Q\partial_w$, this identity is
\[
 \mu wP+\lambda zQ=0.
\]
Putting successively $z=0$ and $w=0$ gives $z\mid P$ and $w\mid Q$, and therefore
\[
 X=h_0(z,w)X_0
\]
with $h_0$ holomorphic. Comparison of linear parts gives $h_0(0)=1$, so $h_0$ is a unit.
\end{proof}

\end{document}
